\begin{figure}[H]
\centering
\begin{tikzpicture}[node distance=8mm]
  \node[moe token] (tok1) at (0,0.45) {};
  \node[moe token] (tok2) at (0,-0.45) {};
  \node[moe router, right=14mm of tok1, yshift=-4.5mm, minimum width=22mm] (router) {Router\\or mechanism};
  \node[moe expert, right=17mm of router, yshift=7mm, minimum width=22mm] (e1) {Expert\\path A};
  \node[moe active expert, right=17mm of router, yshift=-7mm, minimum width=22mm] (e2) {Selected\\path B};
  \node[moe output, right=18mm of router, xshift=29mm, minimum width=22mm] (out) {Output\\check};
  \draw[moe route] (tok1) -- (router);
  \draw[moe route] (tok2) -- (router);
  \draw[moe faint arrow] (router) -- (e1);
  \draw[moe active route] (router) -- (e2);
  \draw[moe arrow] (e1) -- (out);
  \draw[moe arrow] (e2) -- (out);
  \node[align=center, text=BookDeepBlue] at (4.8,-2.0) {Ch. 2.4: Dense feed-forward block as the replacement target};
\end{tikzpicture}
\caption{Dense FFN applied independently to each token row.}
\label{fig:ch02-dense-feed-forward-block-as-the-replacement-target}
\end{figure}
